% ===========================================================================
\chapter{Where to go next}\label{ch:next}

Kestrel Island has 275 nodes. OpenStreetMap has more than nine billion,
and a phone can still plan a route across a continent in a blink. The ideas
in this tutorial are the ones the real thing is built on, but it adds a
great deal on top. Here is a tour of the main differences, each a good
place to dig further, followed by some things to try on the island.

\section{Real map data}

Our map format is a cousin of OpenStreetMap's, and real data can be
converted to it. OpenStreetMap's data can be downloaded for any area, as
XML or in a compact binary format, from its website (the Export tab), from
the Overpass API, or, for whole countries, from sites such as Geofabrik.
A road is a \emph{way} with a \code{highway} tag (\code{primary},
\code{residential}, \code{track}, \code{footway} and many more), a name
tag, and often a \code{oneway=yes} tag and a \code{maxspeed}. Writing a
converter for a small town is a satisfying weekend project, and the
navigator will then guide you round your own streets.

Real data is messier than ours. Roads cross on bridges without sharing a
node, one-way streets are tagged in both directions by mistake, and
footpaths stop a metre short of the road they were meant to join. The
checks of Chapter~\ref{ch:checks}, especially ``every place can reach every
other'', find a surprising number of such problems.

\section{Bigger maps}

\textbf{Projections.} Our local projection is accurate to a few
centimetres across an island, but to a few metres across a county, and it
falls apart across a country. Real systems use projections designed for
the job: Web Mercator for drawing, as in Chapter~\ref{ch:projection}, and
for measuring, either the sphere directly, with the haversine formula, or
a projection made for a region, such as one of the sixty zones of the
Universal Transverse Mercator system, each six degrees of longitude wide.

\textbf{Finding the nearest node.} We check every node. With a billion
nodes, that is far too slow. A \emph{spatial index} divides the map into
squares (or nested rectangles, in an \emph{R-tree}) and keeps a list of
the nodes in each, so a search only looks at the few nearby.

\textbf{Drawing.} A map of the world can't be drawn from scratch every
frame. Web maps cut the world into \emph{tiles}, squares of 256 pixels,
drawn once for each zoom level: one tile for the whole world at zoom 0,
four at zoom 1, and $4^z$ at zoom $z$. At low zoom, most of the detail is
left out, and lines are \emph{simplified}, with fewer points, by
algorithms such as Douglas--Peucker. Modern maps send \emph{vector tiles}
instead, the geometry itself, and the phone draws them, which is why
labels stay upright when you rotate the map.

\section{Faster routes}

Dijkstra's algorithm and A* are where every route planner starts, but on a
continent even A* explores millions of nodes. Real planners do a great deal
of the work in advance.

\begin{itemize}
  \item \textbf{Fewer nodes.} Most nodes are just points along a road,
    with one way in and one way out. Merging each chain of them into a
    single edge shrinks the graph several times over, with no change to
    any route.
  \item \textbf{Searching from both ends.} \emph{Bidirectional} search
    runs one search forwards from the start and another backwards from the
    destination, and stops when they meet. Each covers a circle of half
    the radius, so together they explore about half as much.
  \item \textbf{Better estimates.} The \emph{ALT} technique picks a few
    \emph{landmarks} on the edge of the map and computes, in advance, the
    distance from every node to every landmark. The triangle inequality
    then gives an estimate far closer to the truth than a straight line.
  \item \textbf{Shortcuts.} \emph{Contraction hierarchies} rank the
    nodes by importance and, in advance, add shortcut edges that skip
    over the unimportant ones, so that a long route is found mostly along
    shortcuts between important junctions. A query on a continent then
    explores a few hundred nodes, and takes under a millisecond.
\end{itemize}

\section{Better navigation}

\textbf{Turns.} Our graph has no idea that turning left across traffic
takes longer than going straight on, or that a left turn may be banned at
some junction. Real planners model the cost of each \emph{turn}, from one
edge onto the next, which usually means building a graph whose nodes are
the edges of the road graph.

\textbf{Time.} Our speeds never change. Real ones depend on the time of
day and on live traffic, so the cost of an edge depends on when you reach
it, and the estimated arrival time is updated as the traffic changes.

\textbf{Map matching.} Our navigator matches each fix to the nearest point
of the route. Real ones consider several candidate roads for each fix, and
choose the sequence of roads that best explains all the recent fixes
together, taking into account how far apart they are and how fast the
traveller could have gone between them. The standard method uses a
\emph{hidden Markov model}, and it can follow a car through a city even
when the fixes are tens of metres out.

\textbf{Instructions.} ``Keep left at the fork''; ``take the third exit at
the roundabout''; ``turn right after the bridge''. Good instructions depend
on the shape of each junction and on landmarks, and spoken ones on timing:
the first warning comes early enough to change lanes, the last just
before the turn.

\section{Things to try}

\begin{enumerate}
  \item \textbf{Change the island.} Add a road to \code{kestrel_island.map},
    or make one one-way, and run the checks. (Did you keep every place
    reachable from every other?) Add a place of your own.
  \item \textbf{A ferry.} Add a new kind of road, \code{ferry}, from the
    harbour to Fisher's Cove, slow but direct, and draw it as a dashed
    blue line. When is it quicker than driving round?
  \item \textbf{Avoid tracks.} Add an option to \code{route.nex} that
    makes tracks twice as expensive, without making them slower. Which
    estimate is still safe to use?
  \item \textbf{Bidirectional Dijkstra.} Write it, add it to
    \code{find_route}, and add a check that it agrees with the others.
  \item \textbf{Simplify the coast.} Write the Douglas--Peucker algorithm:
    keep a line's two ends; find the point furthest from the straight line
    between them; if it is further than some tolerance, keep it, and
    simplify the two halves the same way. Use it to draw the coastline
    with fewer points when zoomed out.
  \item \textbf{The frontier.} In the live map, draw the search's frontier
    (the nodes in the queue) as well as the nodes it has explored. You will
    need to record it in the search.
  \item \textbf{Faster but worse.} Chapter~\ref{ch:checks} showed that an
    estimate that overshoots breaks A*'s bookkeeping. Change \code{a_star}
    so that it never updates a node that is already finished, and try the
    doubled estimate again. Now the postcondition holds; which checks fail
    instead? How much less does A* explore, and how much worse are its
    routes? (This trade, a little quality for a lot of speed, is called
    \emph{weighted A*}, and games use it all the time.)
\end{enumerate}
